\chapter*{Abstract}
\addcontentsline{toc}{chapter}{Abstract}

The increasing scale and complexity of modern distributed systems demand robust computational methods for ensuring resilience against partial failures, network partitions, and cascading error propagation. This dissertation develops a unified framework for modeling, analyzing, and optimizing resilience in distributed architectures through three complementary approaches. First, we introduce a graph-theoretic formalism that captures failure propagation dynamics across heterogeneous topologies. Second, we present the \textbf{Nexa Resilience Protocol} (NRP), a lightweight middleware layer that instruments existing distributed systems with adaptive checkpointing and speculative re-execution capabilities. Third, we derive probabilistic bounds on recovery time under varying workload distributions and failure models.

We validate our framework through extensive simulation on synthetic topologies of up to 10,000 nodes and through deployment on a production cluster at the Nimbus Cloud Platform. Our results demonstrate a 42\% reduction in mean time to recovery compared to state-of-the-art Paxos-based approaches, while incurring less than 3\% throughput overhead during normal operation. The techniques developed here generalize across database sharding, microservice orchestration, and edge computing domains, providing practitioners with principled tools for building systems that degrade gracefully under stress.

\vspace{20pt}
\noindent\textbf{Keywords:} distributed systems, resilience, fault tolerance, checkpointing, graph theory, recovery protocols
